% =====================================================================
%  CHƯƠNG 1 — TỔNG QUAN ĐỀ TÀI
%
%  File chương theo cấu trúc template Overleaf (hướng Ứng dụng).
%  Cách dùng: đặt file này trong thư mục Chuong/ của template và thêm
%      % =====================================================================
%  CHƯƠNG 1 — TỔNG QUAN ĐỀ TÀI
%
%  File chương theo cấu trúc template Overleaf (hướng Ứng dụng).
%  Cách dùng: đặt file này trong thư mục Chuong/ của template và thêm
%      % =====================================================================
%  CHƯƠNG 1 — TỔNG QUAN ĐỀ TÀI
%
%  File chương theo cấu trúc template Overleaf (hướng Ứng dụng).
%  Cách dùng: đặt file này trong thư mục Chuong/ của template và thêm
%      % =====================================================================
%  CHƯƠNG 1 — TỔNG QUAN ĐỀ TÀI
%
%  File chương theo cấu trúc template Overleaf (hướng Ứng dụng).
%  Cách dùng: đặt file này trong thư mục Chuong/ của template và thêm
%      \include{Chuong/Chuong1}
%  vào main.tex (sau phần mở đầu, trước các chương sau).
%
%  Các nguồn trích dẫn (\cite) trong chương nằm trong references.bib.
%  KHÔNG hardcode số chương/mục/bảng — luôn dùng \label + \ref / \cite.
% =====================================================================

\chapter{TỔNG QUAN ĐỀ TÀI}
\label{chuong:tong-quan}

Chương này trình bày bối cảnh và lý do hình thành đề tài. Trước hết,
mục~\ref{sec:dat-van-de} đặt ra bài toán học từ vựng tiếng Anh và những
khó khăn cốt lõi của người học. Tiếp theo, mục~\ref{sec:khao-sat-hien-trang}
khảo sát hiện trạng các phương pháp học từ vựng đang được sử dụng, còn
mục~\ref{sec:he-thong-lien-quan} phân tích một số hệ thống tiêu biểu trên
thị trường. Trên cơ sở đó, mục~\ref{sec:phan-tich-nhu-cau} phân tích nhu
cầu thực tế và mục~\ref{sec:muc-tieu} xác định mục tiêu mà hệ thống đề xuất
hướng tới.

% ---------------------------------------------------------------------
\section{Đặt vấn đề}
\label{sec:dat-van-de}

Trong quá trình học một ngoại ngữ, từ vựng đóng vai trò nền tảng. Người
học có thể giao tiếp ở mức cơ bản khi nắm vững từ vựng dù ngữ pháp chưa
hoàn chỉnh, nhưng rất khó diễn đạt ý tưởng nếu vốn từ nghèo nàn. Đối với
tiếng Anh — ngôn ngữ phổ biến nhất trong học thuật, công việc và Internet
— việc tích lũy và ghi nhớ một vốn từ vựng đủ rộng là yêu cầu thiết yếu
đối với học sinh, sinh viên và người đi làm tại Việt Nam.

Tuy nhiên, ghi nhớ từ vựng lâu dài là một thách thức. Theo nghiên cứu kinh
điển của Ebbinghaus về \textit{đường cong lãng quên} (forgetting curve),
phần lớn thông tin mới sẽ bị quên đi nhanh chóng nếu không được ôn tập lại
một cách có hệ thống \cite{ebbinghaus1913memory}. Các nghiên cứu về tâm lý
học nhận thức sau này cũng chỉ ra rằng việc \textit{ôn tập ngắt quãng}
(spaced repetition) — phân bố các lần ôn tập theo những khoảng thời gian
tăng dần thay vì học dồn — giúp cải thiện đáng kể khả năng ghi nhớ dài hạn
so với cách học nhồi nhét \cite{cepeda2006distributed}. Đây chính là cơ sở
khoa học cho các thuật toán lập lịch ôn tập được ứng dụng trong nhiều phần
mềm học tập hiện đại.

Mặc dù vậy, cách học từ vựng phổ biến hiện nay của nhiều người Việt vẫn
còn nhiều hạn chế:

\begin{itemize}
  \item \textbf{Học thụ động và rời rạc:} ghi từ ra sổ tay hoặc tra từ
        điển một lần rồi không có cơ chế nhắc ôn tập, dẫn đến nhanh quên.
  \item \textbf{Không cá nhân hóa:} mọi từ được ôn với tần suất như nhau,
        không phân biệt từ người học đã thuộc với từ còn yếu, gây lãng phí
        thời gian.
  \item \textbf{Thiếu luyện tập chủ động:} người học chủ yếu \emph{nhận
        biết} nghĩa của từ (passive) mà ít có cơ hội \emph{sử dụng} từ
        trong câu hoặc luyện phát âm, trong khi đây mới là kỹ năng quyết
        định khả năng vận dụng thực tế.
  \item \textbf{Khó duy trì động lực:} không có công cụ theo dõi tiến độ,
        chuỗi ngày học hay mục tiêu, khiến người học dễ bỏ cuộc giữa chừng.
\end{itemize}

Từ những vấn đề trên, đề tài hướng tới việc \textbf{xây dựng một hệ thống
ứng dụng học từ vựng tiếng Anh} kết hợp giữa cơ chế ôn tập ngắt quãng có
cơ sở khoa học, khả năng cá nhân hóa nội dung học, các hình thức luyện tập
chủ động (viết câu, luyện phát âm có chấm điểm) và công cụ theo dõi tiến độ,
nhằm giúp người học ghi nhớ từ vựng hiệu quả và bền vững hơn.

% ---------------------------------------------------------------------
\section{Khảo sát hiện trạng}
\label{sec:khao-sat-hien-trang}

Để hiểu rõ bài toán, đề tài khảo sát các phương pháp và công cụ học từ
vựng đang được người học sử dụng phổ biến, cùng những ưu, nhược điểm của
chúng.

\subsection{Các phương pháp học từ vựng phổ biến}

\begin{itemize}
  \item \textbf{Sổ tay và giáo trình giấy:} người học chép từ mới kèm
        nghĩa và ví dụ ra giấy. Cách làm này đơn giản, không phụ thuộc
        thiết bị, nhưng không có cơ chế nhắc ôn tập, khó tra cứu và không
        thống kê được tiến độ.
  \item \textbf{Thẻ ghi nhớ giấy (flashcard):} viết từ ở một mặt và nghĩa
        ở mặt còn lại để tự kiểm tra. Hiệu quả hơn sổ tay nhờ tính chủ
        động, nhưng việc sắp xếp thứ tự ôn tập hoàn toàn thủ công và tốn
        công tạo thẻ.
  \item \textbf{Ứng dụng từ điển:} tra nghĩa, phiên âm và phát âm nhanh
        chóng, nhưng bản chất là công cụ tra cứu tức thời, không được
        thiết kế để hỗ trợ ghi nhớ lâu dài.
  \item \textbf{Ứng dụng học ngoại ngữ trên di động:} cung cấp lộ trình
        học, trò chơi hóa và nhắc nhở. Đây là hướng tiếp cận hiện đại
        nhất và là đối tượng được phân tích kỹ ở mục~\ref{sec:he-thong-lien-quan}.
\end{itemize}

\subsection{Những hạn chế chung}

Qua khảo sát, các phương pháp truyền thống và một số công cụ hiện có còn
tồn tại những hạn chế chung sau:

\begin{enumerate}
  \item Phần lớn dựa trên việc \emph{lặp lại máy móc}, thiếu thuật toán
        lập lịch ôn tập tối ưu theo mức độ ghi nhớ của từng từ.
  \item Nội dung học thường \emph{cố định và không thể tùy biến}; người
        học khó tự bổ sung những từ vựng gắn với nhu cầu riêng (chuyên
        ngành, công việc, sở thích).
  \item Việc tạo nội dung học (định nghĩa, ví dụ, phiên âm, âm thanh) tốn
        nhiều \emph{công sức thủ công}, là rào cản lớn khi muốn mở rộng
        kho từ vựng cá nhân.
  \item Các kỹ năng \emph{sản sinh ngôn ngữ} như đặt câu và phát âm hiếm
        khi được luyện tập kèm phản hồi định lượng để người học biết mình
        đang ở đâu.
  \item Thiếu sự gắn kết giữa \emph{theo dõi tiến độ} và \emph{động lực}:
        người học không có chuỗi ngày học, biểu đồ hoạt động hay yếu tố
        cộng đồng để duy trì thói quen.
\end{enumerate}

Những hạn chế này là cơ sở để đề tài xác định các chức năng cần xây dựng,
được trình bày trong mục~\ref{sec:phan-tich-nhu-cau}.

% ---------------------------------------------------------------------
\section{Các hệ thống liên quan}
\label{sec:he-thong-lien-quan}

Trên thị trường hiện có nhiều ứng dụng hỗ trợ học từ vựng và ngoại ngữ.
Phần này phân tích một số hệ thống tiêu biểu để rút ra ưu, nhược điểm và
khoảng trống mà đề tài có thể khai thác.

\subsection{Duolingo}

Duolingo là một trong những nền tảng học ngoại ngữ miễn phí phổ biến nhất,
nổi bật với cách tiếp cận \textit{trò chơi hóa} (gamification): lộ trình
theo cấp độ, điểm kinh nghiệm, chuỗi ngày học và bảng xếp hạng để duy trì
động lực \cite{vonahn2013duolingo}. Tuy nhiên, Duolingo tập trung vào các
bài học có sẵn theo lộ trình cố định; người học khó tự thêm danh sách từ
vựng riêng và việc đào sâu một từ cụ thể (nhiều nghĩa, ví dụ, ngữ cảnh)
còn hạn chế.

\subsection{Anki}

Anki là phần mềm thẻ ghi nhớ mã nguồn mở, áp dụng thuật toán ôn tập ngắt
quãng SM-2 có nguồn gốc từ dự án SuperMemo \cite{wozniak1990optimization,
anki}. Điểm mạnh của Anki là cơ chế lập lịch ôn tập mạnh mẽ và khả năng
tùy biến thẻ rất cao. Nhược điểm là giao diện và quy trình cấu hình tương
đối phức tạp với người mới; toàn bộ nội dung thẻ phải tự tạo thủ công; và
Anki không tích hợp sẵn cơ chế luyện viết câu hay chấm điểm phát âm.

\subsection{Quizlet}

Quizlet cho phép tạo các bộ thẻ từ vựng và học qua nhiều chế độ (thẻ ghi
nhớ, trắc nghiệm, ghép cặp) cùng tính năng chia sẻ bộ thẻ trong cộng đồng
\cite{quizlet}. Quizlet dễ sử dụng và có kho nội dung do người dùng đóng
góp lớn, nhưng cơ chế ôn tập ngắt quãng không sâu bằng Anki và cũng thiếu
các hình thức luyện tập chủ động có chấm điểm.

\subsection{So sánh và khoảng trống}

Bảng~\ref{tab:so-sanh-he-thong} so sánh các hệ thống trên với hệ thống đề
xuất theo những tiêu chí cốt lõi rút ra từ mục~\ref{sec:khao-sat-hien-trang}.

\begin{table}[h]
  \centering
  \small
  \renewcommand{\arraystretch}{1.3}
  \begin{tabular}{|p{4.1cm}|p{2.0cm}|p{1.7cm}|p{1.9cm}|p{2.3cm}|}
    \hline
    \textbf{Tiêu chí} & \textbf{Duolingo} & \textbf{Anki} & \textbf{Quizlet} & \textbf{Hệ thống đề xuất} \\
    \hline
    Ôn tập ngắt quãng (SRS)          & Hạn chế & Mạnh   & Trung bình & Mạnh \\
    \hline
    Tùy biến / tạo từ vựng cá nhân   & Không   & Có     & Có         & Có \\
    \hline
    Hỗ trợ sinh nội dung bằng AI     & Không   & Không  & Hạn chế    & Có \\
    \hline
    Luyện viết câu có phản hồi       & Không   & Không  & Không      & Có \\
    \hline
    Luyện phát âm có chấm điểm       & Hạn chế & Không  & Không      & Có \\
    \hline
    Theo dõi tiến độ \& gamification & Mạnh    & Cơ bản & Cơ bản     & Có \\
    \hline
  \end{tabular}
  \caption{So sánh các hệ thống học từ vựng tiêu biểu với hệ thống đề xuất}
  \label{tab:so-sanh-he-thong}
\end{table}

Như Bảng~\ref{tab:so-sanh-he-thong} cho thấy, mỗi hệ thống có thế mạnh
riêng nhưng chưa hệ thống nào đồng thời đáp ứng tốt cả bốn yếu tố: kho từ
vựng vừa biên tập sẵn vừa cho phép cá nhân hóa, cơ chế ôn tập ngắt quãng
mạnh, hỗ trợ sinh nội dung tự động bằng AI, và tích hợp luyện viết câu cùng
luyện phát âm có chấm điểm. Đây chính là khoảng trống mà đề tài hướng tới.

% ---------------------------------------------------------------------
\section{Phân tích nhu cầu thực tế}
\label{sec:phan-tich-nhu-cau}

Từ hiện trạng và phân tích các hệ thống liên quan, đề tài xác định các nhu
cầu thực tế cần được hệ thống đáp ứng, chia thành nhu cầu chức năng và nhu
cầu phi chức năng.

\subsection{Đối tượng người dùng}

Hệ thống phục vụ hai nhóm người dùng chính:

\begin{itemize}
  \item \textbf{Người học (user):} người dùng cuối, có nhu cầu học và ôn
        tập từ vựng, tạo từ vựng và bộ thẻ cá nhân, luyện tập và theo dõi
        tiến độ của bản thân.
  \item \textbf{Quản trị viên (admin):} người biên tập và quản lý kho từ
        vựng hệ thống (nội dung được kiểm duyệt, dùng chung cho mọi người
        học), quản lý chủ đề và bộ thẻ mẫu.
\end{itemize}

\subsection{Nhu cầu chức năng}

\begin{enumerate}
  \item \textbf{Quản lý tài khoản và xác thực:} đăng ký, đăng nhập bằng
        email/mật khẩu và qua nhà cung cấp bên thứ ba (Google, Apple,
        GitHub); quản lý hồ sơ và thiết lập mục tiêu học tập.
  \item \textbf{Quản lý từ vựng:} mô hình hóa từ vựng phong phú gồm nhiều
        nghĩa, bản dịch, câu ví dụ, phiên âm, âm thanh và gắn chủ đề; cho
        phép người học tự tạo từ vựng cá nhân bên cạnh kho hệ thống.
  \item \textbf{Tổ chức học liệu:} nhóm từ vựng thành các bộ thẻ (deck)
        và chủ đề; hỗ trợ chia sẻ, sao chép bộ thẻ trong cộng đồng.
  \item \textbf{Học theo cơ chế ôn tập ngắt quãng:} lập lịch ôn tập tự
        động theo mức độ ghi nhớ từng từ, với nhiều dạng câu hỏi luyện tập
        khác nhau trong một phiên học.
  \item \textbf{Luyện tập chủ động có phản hồi:} luyện đặt câu với từ vựng
        và luyện phát âm, đều có cơ chế chấm điểm để người học nhận phản
        hồi định lượng.
  \item \textbf{Hỗ trợ tạo nội dung bằng AI:} sinh nhanh dữ liệu cho một
        từ chỉ từ dạng nguyên thể, nhập hàng loạt và làm giàu dữ liệu từ
        điển, sinh âm thanh phát âm tự động, nhằm giảm công sức tạo nội
        dung thủ công.
  \item \textbf{Theo dõi tiến độ và tạo động lực:} thống kê số từ theo
        trạng thái học, chuỗi ngày học, biểu đồ hoạt động và bảng xếp hạng
        cộng đồng.
\end{enumerate}

\subsection{Nhu cầu phi chức năng}

\begin{itemize}
  \item \textbf{Bảo mật:} bảo vệ tài khoản và dữ liệu người dùng bằng cơ
        chế xác thực bằng token, phân quyền theo vai trò và kiểm soát truy
        cập theo quyền sở hữu dữ liệu.
  \item \textbf{Hiệu năng và khả năng mở rộng:} đáp ứng tốt khi số lượng
        người dùng và lượng từ vựng tăng; các tác vụ nặng (sinh nội dung,
        sinh âm thanh) được xử lý bất đồng bộ để không chặn luồng chính.
  \item \textbf{Tính nhất quán dữ liệu:} đảm bảo toàn vẹn dữ liệu giữa từ
        vựng, bộ thẻ và tiến độ học của người dùng.
  \item \textbf{Khả năng bảo trì và mở rộng:} kiến trúc rõ ràng, API có
        phiên bản hóa để dễ dàng phát triển thêm tính năng về sau.
\end{itemize}

% ---------------------------------------------------------------------
\section{Mục tiêu của hệ thống}
\label{sec:muc-tieu}

\subsection{Mục tiêu tổng quát}

Mục tiêu tổng quát của đề tài là \textbf{xây dựng hệ thống phía máy chủ
(backend) cho ứng dụng học từ vựng tiếng Anh}, cung cấp đầy đủ các dịch vụ
qua giao diện lập trình ứng dụng (API) để kết hợp giữa kho từ vựng được
biên tập, khả năng cá nhân hóa, cơ chế ôn tập ngắt quãng có cơ sở khoa
học, hỗ trợ sinh nội dung bằng AI và các hình thức luyện tập chủ động có
chấm điểm, qua đó giúp người học ghi nhớ và vận dụng từ vựng hiệu quả hơn.

\subsection{Mục tiêu cụ thể}

\begin{enumerate}
  \item Xây dựng phân hệ \textbf{tài khoản và xác thực} hỗ trợ nhiều
        phương thức đăng nhập, phân quyền người dùng và quản trị viên.
  \item Thiết kế \textbf{mô hình dữ liệu từ vựng} đầy đủ (nhiều nghĩa, bản
        dịch, ví dụ, phiên âm, âm thanh, chủ đề) cho cả kho hệ thống và từ
        vựng cá nhân.
  \item Xây dựng phân hệ \textbf{bộ thẻ và chủ đề}, hỗ trợ chia sẻ và sao
        chép trong cộng đồng.
  \item Cài đặt \textbf{động cơ học theo ôn tập ngắt quãng} dựa trên thuật
        toán SM-2 kết hợp các bước học (learning steps), với nhiều dạng
        câu hỏi luyện tập trong một phiên học.
  \item Xây dựng các phân hệ \textbf{luyện viết câu} và \textbf{luyện phát
        âm} có cơ chế chấm điểm và phản hồi.
  \item Tích hợp các tính năng \textbf{hỗ trợ bằng AI}: tạo nhanh và nhập
        hàng loạt từ vựng, làm giàu dữ liệu từ điển và sinh âm thanh phát
        âm tự động.
  \item Xây dựng phân hệ \textbf{theo dõi tiến độ}: thống kê, chuỗi ngày
        học, biểu đồ hoạt động và bảng xếp hạng.
  \item Đảm bảo các \textbf{yêu cầu phi chức năng} về bảo mật, hiệu năng,
        khả năng mở rộng và tính nhất quán dữ liệu.
\end{enumerate}

\subsection{Phạm vi đề tài}

Đề tài bao gồm cả \textbf{hệ thống phía máy chủ (backend)} cùng các API
phục vụ toàn bộ nghiệp vụ nêu trên và \textbf{ứng dụng giao diện người dùng
phía máy khách (frontend)} tiêu thụ các API đó. Trong đó, trọng tâm trình
bày đặt ở phần backend --- nơi tập trung phần lớn logic nghiệp vụ, mô hình
dữ liệu và các tính năng thông minh; phần frontend được trình bày ở mức kiến
trúc và hiện thực hóa tầng giao diện, đủ để làm rõ cách toàn bộ hệ thống vận
hành từ người dùng cuối đến cơ sở dữ liệu. Các chương tiếp theo sẽ lần lượt
trình bày cơ sở lý thuyết, phân tích và thiết kế hệ thống, cũng như quá
trình hiện thực hóa và kiểm thử.

%  vào main.tex (sau phần mở đầu, trước các chương sau).
%
%  Các nguồn trích dẫn (\cite) trong chương nằm trong references.bib.
%  KHÔNG hardcode số chương/mục/bảng — luôn dùng \label + \ref / \cite.
% =====================================================================

\chapter{TỔNG QUAN ĐỀ TÀI}
\label{chuong:tong-quan}

Chương này trình bày bối cảnh và lý do hình thành đề tài. Trước hết,
mục~\ref{sec:dat-van-de} đặt ra bài toán học từ vựng tiếng Anh và những
khó khăn cốt lõi của người học. Tiếp theo, mục~\ref{sec:khao-sat-hien-trang}
khảo sát hiện trạng các phương pháp học từ vựng đang được sử dụng, còn
mục~\ref{sec:he-thong-lien-quan} phân tích một số hệ thống tiêu biểu trên
thị trường. Trên cơ sở đó, mục~\ref{sec:phan-tich-nhu-cau} phân tích nhu
cầu thực tế và mục~\ref{sec:muc-tieu} xác định mục tiêu mà hệ thống đề xuất
hướng tới.

% ---------------------------------------------------------------------
\section{Đặt vấn đề}
\label{sec:dat-van-de}

Trong quá trình học một ngoại ngữ, từ vựng đóng vai trò nền tảng. Người
học có thể giao tiếp ở mức cơ bản khi nắm vững từ vựng dù ngữ pháp chưa
hoàn chỉnh, nhưng rất khó diễn đạt ý tưởng nếu vốn từ nghèo nàn. Đối với
tiếng Anh — ngôn ngữ phổ biến nhất trong học thuật, công việc và Internet
— việc tích lũy và ghi nhớ một vốn từ vựng đủ rộng là yêu cầu thiết yếu
đối với học sinh, sinh viên và người đi làm tại Việt Nam.

Tuy nhiên, ghi nhớ từ vựng lâu dài là một thách thức. Theo nghiên cứu kinh
điển của Ebbinghaus về \textit{đường cong lãng quên} (forgetting curve),
phần lớn thông tin mới sẽ bị quên đi nhanh chóng nếu không được ôn tập lại
một cách có hệ thống \cite{ebbinghaus1913memory}. Các nghiên cứu về tâm lý
học nhận thức sau này cũng chỉ ra rằng việc \textit{ôn tập ngắt quãng}
(spaced repetition) — phân bố các lần ôn tập theo những khoảng thời gian
tăng dần thay vì học dồn — giúp cải thiện đáng kể khả năng ghi nhớ dài hạn
so với cách học nhồi nhét \cite{cepeda2006distributed}. Đây chính là cơ sở
khoa học cho các thuật toán lập lịch ôn tập được ứng dụng trong nhiều phần
mềm học tập hiện đại.

Mặc dù vậy, cách học từ vựng phổ biến hiện nay của nhiều người Việt vẫn
còn nhiều hạn chế:

\begin{itemize}
  \item \textbf{Học thụ động và rời rạc:} ghi từ ra sổ tay hoặc tra từ
        điển một lần rồi không có cơ chế nhắc ôn tập, dẫn đến nhanh quên.
  \item \textbf{Không cá nhân hóa:} mọi từ được ôn với tần suất như nhau,
        không phân biệt từ người học đã thuộc với từ còn yếu, gây lãng phí
        thời gian.
  \item \textbf{Thiếu luyện tập chủ động:} người học chủ yếu \emph{nhận
        biết} nghĩa của từ (passive) mà ít có cơ hội \emph{sử dụng} từ
        trong câu hoặc luyện phát âm, trong khi đây mới là kỹ năng quyết
        định khả năng vận dụng thực tế.
  \item \textbf{Khó duy trì động lực:} không có công cụ theo dõi tiến độ,
        chuỗi ngày học hay mục tiêu, khiến người học dễ bỏ cuộc giữa chừng.
\end{itemize}

Từ những vấn đề trên, đề tài hướng tới việc \textbf{xây dựng một hệ thống
ứng dụng học từ vựng tiếng Anh} kết hợp giữa cơ chế ôn tập ngắt quãng có
cơ sở khoa học, khả năng cá nhân hóa nội dung học, các hình thức luyện tập
chủ động (viết câu, luyện phát âm có chấm điểm) và công cụ theo dõi tiến độ,
nhằm giúp người học ghi nhớ từ vựng hiệu quả và bền vững hơn.

% ---------------------------------------------------------------------
\section{Khảo sát hiện trạng}
\label{sec:khao-sat-hien-trang}

Để hiểu rõ bài toán, đề tài khảo sát các phương pháp và công cụ học từ
vựng đang được người học sử dụng phổ biến, cùng những ưu, nhược điểm của
chúng.

\subsection{Các phương pháp học từ vựng phổ biến}

\begin{itemize}
  \item \textbf{Sổ tay và giáo trình giấy:} người học chép từ mới kèm
        nghĩa và ví dụ ra giấy. Cách làm này đơn giản, không phụ thuộc
        thiết bị, nhưng không có cơ chế nhắc ôn tập, khó tra cứu và không
        thống kê được tiến độ.
  \item \textbf{Thẻ ghi nhớ giấy (flashcard):} viết từ ở một mặt và nghĩa
        ở mặt còn lại để tự kiểm tra. Hiệu quả hơn sổ tay nhờ tính chủ
        động, nhưng việc sắp xếp thứ tự ôn tập hoàn toàn thủ công và tốn
        công tạo thẻ.
  \item \textbf{Ứng dụng từ điển:} tra nghĩa, phiên âm và phát âm nhanh
        chóng, nhưng bản chất là công cụ tra cứu tức thời, không được
        thiết kế để hỗ trợ ghi nhớ lâu dài.
  \item \textbf{Ứng dụng học ngoại ngữ trên di động:} cung cấp lộ trình
        học, trò chơi hóa và nhắc nhở. Đây là hướng tiếp cận hiện đại
        nhất và là đối tượng được phân tích kỹ ở mục~\ref{sec:he-thong-lien-quan}.
\end{itemize}

\subsection{Những hạn chế chung}

Qua khảo sát, các phương pháp truyền thống và một số công cụ hiện có còn
tồn tại những hạn chế chung sau:

\begin{enumerate}
  \item Phần lớn dựa trên việc \emph{lặp lại máy móc}, thiếu thuật toán
        lập lịch ôn tập tối ưu theo mức độ ghi nhớ của từng từ.
  \item Nội dung học thường \emph{cố định và không thể tùy biến}; người
        học khó tự bổ sung những từ vựng gắn với nhu cầu riêng (chuyên
        ngành, công việc, sở thích).
  \item Việc tạo nội dung học (định nghĩa, ví dụ, phiên âm, âm thanh) tốn
        nhiều \emph{công sức thủ công}, là rào cản lớn khi muốn mở rộng
        kho từ vựng cá nhân.
  \item Các kỹ năng \emph{sản sinh ngôn ngữ} như đặt câu và phát âm hiếm
        khi được luyện tập kèm phản hồi định lượng để người học biết mình
        đang ở đâu.
  \item Thiếu sự gắn kết giữa \emph{theo dõi tiến độ} và \emph{động lực}:
        người học không có chuỗi ngày học, biểu đồ hoạt động hay yếu tố
        cộng đồng để duy trì thói quen.
\end{enumerate}

Những hạn chế này là cơ sở để đề tài xác định các chức năng cần xây dựng,
được trình bày trong mục~\ref{sec:phan-tich-nhu-cau}.

% ---------------------------------------------------------------------
\section{Các hệ thống liên quan}
\label{sec:he-thong-lien-quan}

Trên thị trường hiện có nhiều ứng dụng hỗ trợ học từ vựng và ngoại ngữ.
Phần này phân tích một số hệ thống tiêu biểu để rút ra ưu, nhược điểm và
khoảng trống mà đề tài có thể khai thác.

\subsection{Duolingo}

Duolingo là một trong những nền tảng học ngoại ngữ miễn phí phổ biến nhất,
nổi bật với cách tiếp cận \textit{trò chơi hóa} (gamification): lộ trình
theo cấp độ, điểm kinh nghiệm, chuỗi ngày học và bảng xếp hạng để duy trì
động lực \cite{vonahn2013duolingo}. Tuy nhiên, Duolingo tập trung vào các
bài học có sẵn theo lộ trình cố định; người học khó tự thêm danh sách từ
vựng riêng và việc đào sâu một từ cụ thể (nhiều nghĩa, ví dụ, ngữ cảnh)
còn hạn chế.

\subsection{Anki}

Anki là phần mềm thẻ ghi nhớ mã nguồn mở, áp dụng thuật toán ôn tập ngắt
quãng SM-2 có nguồn gốc từ dự án SuperMemo \cite{wozniak1990optimization,
anki}. Điểm mạnh của Anki là cơ chế lập lịch ôn tập mạnh mẽ và khả năng
tùy biến thẻ rất cao. Nhược điểm là giao diện và quy trình cấu hình tương
đối phức tạp với người mới; toàn bộ nội dung thẻ phải tự tạo thủ công; và
Anki không tích hợp sẵn cơ chế luyện viết câu hay chấm điểm phát âm.

\subsection{Quizlet}

Quizlet cho phép tạo các bộ thẻ từ vựng và học qua nhiều chế độ (thẻ ghi
nhớ, trắc nghiệm, ghép cặp) cùng tính năng chia sẻ bộ thẻ trong cộng đồng
\cite{quizlet}. Quizlet dễ sử dụng và có kho nội dung do người dùng đóng
góp lớn, nhưng cơ chế ôn tập ngắt quãng không sâu bằng Anki và cũng thiếu
các hình thức luyện tập chủ động có chấm điểm.

\subsection{So sánh và khoảng trống}

Bảng~\ref{tab:so-sanh-he-thong} so sánh các hệ thống trên với hệ thống đề
xuất theo những tiêu chí cốt lõi rút ra từ mục~\ref{sec:khao-sat-hien-trang}.

\begin{table}[h]
  \centering
  \small
  \renewcommand{\arraystretch}{1.3}
  \begin{tabular}{|p{4.1cm}|p{2.0cm}|p{1.7cm}|p{1.9cm}|p{2.3cm}|}
    \hline
    \textbf{Tiêu chí} & \textbf{Duolingo} & \textbf{Anki} & \textbf{Quizlet} & \textbf{Hệ thống đề xuất} \\
    \hline
    Ôn tập ngắt quãng (SRS)          & Hạn chế & Mạnh   & Trung bình & Mạnh \\
    \hline
    Tùy biến / tạo từ vựng cá nhân   & Không   & Có     & Có         & Có \\
    \hline
    Hỗ trợ sinh nội dung bằng AI     & Không   & Không  & Hạn chế    & Có \\
    \hline
    Luyện viết câu có phản hồi       & Không   & Không  & Không      & Có \\
    \hline
    Luyện phát âm có chấm điểm       & Hạn chế & Không  & Không      & Có \\
    \hline
    Theo dõi tiến độ \& gamification & Mạnh    & Cơ bản & Cơ bản     & Có \\
    \hline
  \end{tabular}
  \caption{So sánh các hệ thống học từ vựng tiêu biểu với hệ thống đề xuất}
  \label{tab:so-sanh-he-thong}
\end{table}

Như Bảng~\ref{tab:so-sanh-he-thong} cho thấy, mỗi hệ thống có thế mạnh
riêng nhưng chưa hệ thống nào đồng thời đáp ứng tốt cả bốn yếu tố: kho từ
vựng vừa biên tập sẵn vừa cho phép cá nhân hóa, cơ chế ôn tập ngắt quãng
mạnh, hỗ trợ sinh nội dung tự động bằng AI, và tích hợp luyện viết câu cùng
luyện phát âm có chấm điểm. Đây chính là khoảng trống mà đề tài hướng tới.

% ---------------------------------------------------------------------
\section{Phân tích nhu cầu thực tế}
\label{sec:phan-tich-nhu-cau}

Từ hiện trạng và phân tích các hệ thống liên quan, đề tài xác định các nhu
cầu thực tế cần được hệ thống đáp ứng, chia thành nhu cầu chức năng và nhu
cầu phi chức năng.

\subsection{Đối tượng người dùng}

Hệ thống phục vụ hai nhóm người dùng chính:

\begin{itemize}
  \item \textbf{Người học (user):} người dùng cuối, có nhu cầu học và ôn
        tập từ vựng, tạo từ vựng và bộ thẻ cá nhân, luyện tập và theo dõi
        tiến độ của bản thân.
  \item \textbf{Quản trị viên (admin):} người biên tập và quản lý kho từ
        vựng hệ thống (nội dung được kiểm duyệt, dùng chung cho mọi người
        học), quản lý chủ đề và bộ thẻ mẫu.
\end{itemize}

\subsection{Nhu cầu chức năng}

\begin{enumerate}
  \item \textbf{Quản lý tài khoản và xác thực:} đăng ký, đăng nhập bằng
        email/mật khẩu và qua nhà cung cấp bên thứ ba (Google, Apple,
        GitHub); quản lý hồ sơ và thiết lập mục tiêu học tập.
  \item \textbf{Quản lý từ vựng:} mô hình hóa từ vựng phong phú gồm nhiều
        nghĩa, bản dịch, câu ví dụ, phiên âm, âm thanh và gắn chủ đề; cho
        phép người học tự tạo từ vựng cá nhân bên cạnh kho hệ thống.
  \item \textbf{Tổ chức học liệu:} nhóm từ vựng thành các bộ thẻ (deck)
        và chủ đề; hỗ trợ chia sẻ, sao chép bộ thẻ trong cộng đồng.
  \item \textbf{Học theo cơ chế ôn tập ngắt quãng:} lập lịch ôn tập tự
        động theo mức độ ghi nhớ từng từ, với nhiều dạng câu hỏi luyện tập
        khác nhau trong một phiên học.
  \item \textbf{Luyện tập chủ động có phản hồi:} luyện đặt câu với từ vựng
        và luyện phát âm, đều có cơ chế chấm điểm để người học nhận phản
        hồi định lượng.
  \item \textbf{Hỗ trợ tạo nội dung bằng AI:} sinh nhanh dữ liệu cho một
        từ chỉ từ dạng nguyên thể, nhập hàng loạt và làm giàu dữ liệu từ
        điển, sinh âm thanh phát âm tự động, nhằm giảm công sức tạo nội
        dung thủ công.
  \item \textbf{Theo dõi tiến độ và tạo động lực:} thống kê số từ theo
        trạng thái học, chuỗi ngày học, biểu đồ hoạt động và bảng xếp hạng
        cộng đồng.
\end{enumerate}

\subsection{Nhu cầu phi chức năng}

\begin{itemize}
  \item \textbf{Bảo mật:} bảo vệ tài khoản và dữ liệu người dùng bằng cơ
        chế xác thực bằng token, phân quyền theo vai trò và kiểm soát truy
        cập theo quyền sở hữu dữ liệu.
  \item \textbf{Hiệu năng và khả năng mở rộng:} đáp ứng tốt khi số lượng
        người dùng và lượng từ vựng tăng; các tác vụ nặng (sinh nội dung,
        sinh âm thanh) được xử lý bất đồng bộ để không chặn luồng chính.
  \item \textbf{Tính nhất quán dữ liệu:} đảm bảo toàn vẹn dữ liệu giữa từ
        vựng, bộ thẻ và tiến độ học của người dùng.
  \item \textbf{Khả năng bảo trì và mở rộng:} kiến trúc rõ ràng, API có
        phiên bản hóa để dễ dàng phát triển thêm tính năng về sau.
\end{itemize}

% ---------------------------------------------------------------------
\section{Mục tiêu của hệ thống}
\label{sec:muc-tieu}

\subsection{Mục tiêu tổng quát}

Mục tiêu tổng quát của đề tài là \textbf{xây dựng hệ thống phía máy chủ
(backend) cho ứng dụng học từ vựng tiếng Anh}, cung cấp đầy đủ các dịch vụ
qua giao diện lập trình ứng dụng (API) để kết hợp giữa kho từ vựng được
biên tập, khả năng cá nhân hóa, cơ chế ôn tập ngắt quãng có cơ sở khoa
học, hỗ trợ sinh nội dung bằng AI và các hình thức luyện tập chủ động có
chấm điểm, qua đó giúp người học ghi nhớ và vận dụng từ vựng hiệu quả hơn.

\subsection{Mục tiêu cụ thể}

\begin{enumerate}
  \item Xây dựng phân hệ \textbf{tài khoản và xác thực} hỗ trợ nhiều
        phương thức đăng nhập, phân quyền người dùng và quản trị viên.
  \item Thiết kế \textbf{mô hình dữ liệu từ vựng} đầy đủ (nhiều nghĩa, bản
        dịch, ví dụ, phiên âm, âm thanh, chủ đề) cho cả kho hệ thống và từ
        vựng cá nhân.
  \item Xây dựng phân hệ \textbf{bộ thẻ và chủ đề}, hỗ trợ chia sẻ và sao
        chép trong cộng đồng.
  \item Cài đặt \textbf{động cơ học theo ôn tập ngắt quãng} dựa trên thuật
        toán SM-2 kết hợp các bước học (learning steps), với nhiều dạng
        câu hỏi luyện tập trong một phiên học.
  \item Xây dựng các phân hệ \textbf{luyện viết câu} và \textbf{luyện phát
        âm} có cơ chế chấm điểm và phản hồi.
  \item Tích hợp các tính năng \textbf{hỗ trợ bằng AI}: tạo nhanh và nhập
        hàng loạt từ vựng, làm giàu dữ liệu từ điển và sinh âm thanh phát
        âm tự động.
  \item Xây dựng phân hệ \textbf{theo dõi tiến độ}: thống kê, chuỗi ngày
        học, biểu đồ hoạt động và bảng xếp hạng.
  \item Đảm bảo các \textbf{yêu cầu phi chức năng} về bảo mật, hiệu năng,
        khả năng mở rộng và tính nhất quán dữ liệu.
\end{enumerate}

\subsection{Phạm vi đề tài}

Đề tài bao gồm cả \textbf{hệ thống phía máy chủ (backend)} cùng các API
phục vụ toàn bộ nghiệp vụ nêu trên và \textbf{ứng dụng giao diện người dùng
phía máy khách (frontend)} tiêu thụ các API đó. Trong đó, trọng tâm trình
bày đặt ở phần backend --- nơi tập trung phần lớn logic nghiệp vụ, mô hình
dữ liệu và các tính năng thông minh; phần frontend được trình bày ở mức kiến
trúc và hiện thực hóa tầng giao diện, đủ để làm rõ cách toàn bộ hệ thống vận
hành từ người dùng cuối đến cơ sở dữ liệu. Các chương tiếp theo sẽ lần lượt
trình bày cơ sở lý thuyết, phân tích và thiết kế hệ thống, cũng như quá
trình hiện thực hóa và kiểm thử.

%  vào main.tex (sau phần mở đầu, trước các chương sau).
%
%  Các nguồn trích dẫn (\cite) trong chương nằm trong references.bib.
%  KHÔNG hardcode số chương/mục/bảng — luôn dùng \label + \ref / \cite.
% =====================================================================

\chapter{TỔNG QUAN ĐỀ TÀI}
\label{chuong:tong-quan}

Chương này trình bày bối cảnh và lý do hình thành đề tài. Trước hết,
mục~\ref{sec:dat-van-de} đặt ra bài toán học từ vựng tiếng Anh và những
khó khăn cốt lõi của người học. Tiếp theo, mục~\ref{sec:khao-sat-hien-trang}
khảo sát hiện trạng các phương pháp học từ vựng đang được sử dụng, còn
mục~\ref{sec:he-thong-lien-quan} phân tích một số hệ thống tiêu biểu trên
thị trường. Trên cơ sở đó, mục~\ref{sec:phan-tich-nhu-cau} phân tích nhu
cầu thực tế và mục~\ref{sec:muc-tieu} xác định mục tiêu mà hệ thống đề xuất
hướng tới.

% ---------------------------------------------------------------------
\section{Đặt vấn đề}
\label{sec:dat-van-de}

Trong quá trình học một ngoại ngữ, từ vựng đóng vai trò nền tảng. Người
học có thể giao tiếp ở mức cơ bản khi nắm vững từ vựng dù ngữ pháp chưa
hoàn chỉnh, nhưng rất khó diễn đạt ý tưởng nếu vốn từ nghèo nàn. Đối với
tiếng Anh — ngôn ngữ phổ biến nhất trong học thuật, công việc và Internet
— việc tích lũy và ghi nhớ một vốn từ vựng đủ rộng là yêu cầu thiết yếu
đối với học sinh, sinh viên và người đi làm tại Việt Nam.

Tuy nhiên, ghi nhớ từ vựng lâu dài là một thách thức. Theo nghiên cứu kinh
điển của Ebbinghaus về \textit{đường cong lãng quên} (forgetting curve),
phần lớn thông tin mới sẽ bị quên đi nhanh chóng nếu không được ôn tập lại
một cách có hệ thống \cite{ebbinghaus1913memory}. Các nghiên cứu về tâm lý
học nhận thức sau này cũng chỉ ra rằng việc \textit{ôn tập ngắt quãng}
(spaced repetition) — phân bố các lần ôn tập theo những khoảng thời gian
tăng dần thay vì học dồn — giúp cải thiện đáng kể khả năng ghi nhớ dài hạn
so với cách học nhồi nhét \cite{cepeda2006distributed}. Đây chính là cơ sở
khoa học cho các thuật toán lập lịch ôn tập được ứng dụng trong nhiều phần
mềm học tập hiện đại.

Mặc dù vậy, cách học từ vựng phổ biến hiện nay của nhiều người Việt vẫn
còn nhiều hạn chế:

\begin{itemize}
  \item \textbf{Học thụ động và rời rạc:} ghi từ ra sổ tay hoặc tra từ
        điển một lần rồi không có cơ chế nhắc ôn tập, dẫn đến nhanh quên.
  \item \textbf{Không cá nhân hóa:} mọi từ được ôn với tần suất như nhau,
        không phân biệt từ người học đã thuộc với từ còn yếu, gây lãng phí
        thời gian.
  \item \textbf{Thiếu luyện tập chủ động:} người học chủ yếu \emph{nhận
        biết} nghĩa của từ (passive) mà ít có cơ hội \emph{sử dụng} từ
        trong câu hoặc luyện phát âm, trong khi đây mới là kỹ năng quyết
        định khả năng vận dụng thực tế.
  \item \textbf{Khó duy trì động lực:} không có công cụ theo dõi tiến độ,
        chuỗi ngày học hay mục tiêu, khiến người học dễ bỏ cuộc giữa chừng.
\end{itemize}

Từ những vấn đề trên, đề tài hướng tới việc \textbf{xây dựng một hệ thống
ứng dụng học từ vựng tiếng Anh} kết hợp giữa cơ chế ôn tập ngắt quãng có
cơ sở khoa học, khả năng cá nhân hóa nội dung học, các hình thức luyện tập
chủ động (viết câu, luyện phát âm có chấm điểm) và công cụ theo dõi tiến độ,
nhằm giúp người học ghi nhớ từ vựng hiệu quả và bền vững hơn.

% ---------------------------------------------------------------------
\section{Khảo sát hiện trạng}
\label{sec:khao-sat-hien-trang}

Để hiểu rõ bài toán, đề tài khảo sát các phương pháp và công cụ học từ
vựng đang được người học sử dụng phổ biến, cùng những ưu, nhược điểm của
chúng.

\subsection{Các phương pháp học từ vựng phổ biến}

\begin{itemize}
  \item \textbf{Sổ tay và giáo trình giấy:} người học chép từ mới kèm
        nghĩa và ví dụ ra giấy. Cách làm này đơn giản, không phụ thuộc
        thiết bị, nhưng không có cơ chế nhắc ôn tập, khó tra cứu và không
        thống kê được tiến độ.
  \item \textbf{Thẻ ghi nhớ giấy (flashcard):} viết từ ở một mặt và nghĩa
        ở mặt còn lại để tự kiểm tra. Hiệu quả hơn sổ tay nhờ tính chủ
        động, nhưng việc sắp xếp thứ tự ôn tập hoàn toàn thủ công và tốn
        công tạo thẻ.
  \item \textbf{Ứng dụng từ điển:} tra nghĩa, phiên âm và phát âm nhanh
        chóng, nhưng bản chất là công cụ tra cứu tức thời, không được
        thiết kế để hỗ trợ ghi nhớ lâu dài.
  \item \textbf{Ứng dụng học ngoại ngữ trên di động:} cung cấp lộ trình
        học, trò chơi hóa và nhắc nhở. Đây là hướng tiếp cận hiện đại
        nhất và là đối tượng được phân tích kỹ ở mục~\ref{sec:he-thong-lien-quan}.
\end{itemize}

\subsection{Những hạn chế chung}

Qua khảo sát, các phương pháp truyền thống và một số công cụ hiện có còn
tồn tại những hạn chế chung sau:

\begin{enumerate}
  \item Phần lớn dựa trên việc \emph{lặp lại máy móc}, thiếu thuật toán
        lập lịch ôn tập tối ưu theo mức độ ghi nhớ của từng từ.
  \item Nội dung học thường \emph{cố định và không thể tùy biến}; người
        học khó tự bổ sung những từ vựng gắn với nhu cầu riêng (chuyên
        ngành, công việc, sở thích).
  \item Việc tạo nội dung học (định nghĩa, ví dụ, phiên âm, âm thanh) tốn
        nhiều \emph{công sức thủ công}, là rào cản lớn khi muốn mở rộng
        kho từ vựng cá nhân.
  \item Các kỹ năng \emph{sản sinh ngôn ngữ} như đặt câu và phát âm hiếm
        khi được luyện tập kèm phản hồi định lượng để người học biết mình
        đang ở đâu.
  \item Thiếu sự gắn kết giữa \emph{theo dõi tiến độ} và \emph{động lực}:
        người học không có chuỗi ngày học, biểu đồ hoạt động hay yếu tố
        cộng đồng để duy trì thói quen.
\end{enumerate}

Những hạn chế này là cơ sở để đề tài xác định các chức năng cần xây dựng,
được trình bày trong mục~\ref{sec:phan-tich-nhu-cau}.

% ---------------------------------------------------------------------
\section{Các hệ thống liên quan}
\label{sec:he-thong-lien-quan}

Trên thị trường hiện có nhiều ứng dụng hỗ trợ học từ vựng và ngoại ngữ.
Phần này phân tích một số hệ thống tiêu biểu để rút ra ưu, nhược điểm và
khoảng trống mà đề tài có thể khai thác.

\subsection{Duolingo}

Duolingo là một trong những nền tảng học ngoại ngữ miễn phí phổ biến nhất,
nổi bật với cách tiếp cận \textit{trò chơi hóa} (gamification): lộ trình
theo cấp độ, điểm kinh nghiệm, chuỗi ngày học và bảng xếp hạng để duy trì
động lực \cite{vonahn2013duolingo}. Tuy nhiên, Duolingo tập trung vào các
bài học có sẵn theo lộ trình cố định; người học khó tự thêm danh sách từ
vựng riêng và việc đào sâu một từ cụ thể (nhiều nghĩa, ví dụ, ngữ cảnh)
còn hạn chế.

\subsection{Anki}

Anki là phần mềm thẻ ghi nhớ mã nguồn mở, áp dụng thuật toán ôn tập ngắt
quãng SM-2 có nguồn gốc từ dự án SuperMemo \cite{wozniak1990optimization,
anki}. Điểm mạnh của Anki là cơ chế lập lịch ôn tập mạnh mẽ và khả năng
tùy biến thẻ rất cao. Nhược điểm là giao diện và quy trình cấu hình tương
đối phức tạp với người mới; toàn bộ nội dung thẻ phải tự tạo thủ công; và
Anki không tích hợp sẵn cơ chế luyện viết câu hay chấm điểm phát âm.

\subsection{Quizlet}

Quizlet cho phép tạo các bộ thẻ từ vựng và học qua nhiều chế độ (thẻ ghi
nhớ, trắc nghiệm, ghép cặp) cùng tính năng chia sẻ bộ thẻ trong cộng đồng
\cite{quizlet}. Quizlet dễ sử dụng và có kho nội dung do người dùng đóng
góp lớn, nhưng cơ chế ôn tập ngắt quãng không sâu bằng Anki và cũng thiếu
các hình thức luyện tập chủ động có chấm điểm.

\subsection{So sánh và khoảng trống}

Bảng~\ref{tab:so-sanh-he-thong} so sánh các hệ thống trên với hệ thống đề
xuất theo những tiêu chí cốt lõi rút ra từ mục~\ref{sec:khao-sat-hien-trang}.

\begin{table}[h]
  \centering
  \small
  \renewcommand{\arraystretch}{1.3}
  \begin{tabular}{|p{4.1cm}|p{2.0cm}|p{1.7cm}|p{1.9cm}|p{2.3cm}|}
    \hline
    \textbf{Tiêu chí} & \textbf{Duolingo} & \textbf{Anki} & \textbf{Quizlet} & \textbf{Hệ thống đề xuất} \\
    \hline
    Ôn tập ngắt quãng (SRS)          & Hạn chế & Mạnh   & Trung bình & Mạnh \\
    \hline
    Tùy biến / tạo từ vựng cá nhân   & Không   & Có     & Có         & Có \\
    \hline
    Hỗ trợ sinh nội dung bằng AI     & Không   & Không  & Hạn chế    & Có \\
    \hline
    Luyện viết câu có phản hồi       & Không   & Không  & Không      & Có \\
    \hline
    Luyện phát âm có chấm điểm       & Hạn chế & Không  & Không      & Có \\
    \hline
    Theo dõi tiến độ \& gamification & Mạnh    & Cơ bản & Cơ bản     & Có \\
    \hline
  \end{tabular}
  \caption{So sánh các hệ thống học từ vựng tiêu biểu với hệ thống đề xuất}
  \label{tab:so-sanh-he-thong}
\end{table}

Như Bảng~\ref{tab:so-sanh-he-thong} cho thấy, mỗi hệ thống có thế mạnh
riêng nhưng chưa hệ thống nào đồng thời đáp ứng tốt cả bốn yếu tố: kho từ
vựng vừa biên tập sẵn vừa cho phép cá nhân hóa, cơ chế ôn tập ngắt quãng
mạnh, hỗ trợ sinh nội dung tự động bằng AI, và tích hợp luyện viết câu cùng
luyện phát âm có chấm điểm. Đây chính là khoảng trống mà đề tài hướng tới.

% ---------------------------------------------------------------------
\section{Phân tích nhu cầu thực tế}
\label{sec:phan-tich-nhu-cau}

Từ hiện trạng và phân tích các hệ thống liên quan, đề tài xác định các nhu
cầu thực tế cần được hệ thống đáp ứng, chia thành nhu cầu chức năng và nhu
cầu phi chức năng.

\subsection{Đối tượng người dùng}

Hệ thống phục vụ hai nhóm người dùng chính:

\begin{itemize}
  \item \textbf{Người học (user):} người dùng cuối, có nhu cầu học và ôn
        tập từ vựng, tạo từ vựng và bộ thẻ cá nhân, luyện tập và theo dõi
        tiến độ của bản thân.
  \item \textbf{Quản trị viên (admin):} người biên tập và quản lý kho từ
        vựng hệ thống (nội dung được kiểm duyệt, dùng chung cho mọi người
        học), quản lý chủ đề và bộ thẻ mẫu.
\end{itemize}

\subsection{Nhu cầu chức năng}

\begin{enumerate}
  \item \textbf{Quản lý tài khoản và xác thực:} đăng ký, đăng nhập bằng
        email/mật khẩu và qua nhà cung cấp bên thứ ba (Google, Apple,
        GitHub); quản lý hồ sơ và thiết lập mục tiêu học tập.
  \item \textbf{Quản lý từ vựng:} mô hình hóa từ vựng phong phú gồm nhiều
        nghĩa, bản dịch, câu ví dụ, phiên âm, âm thanh và gắn chủ đề; cho
        phép người học tự tạo từ vựng cá nhân bên cạnh kho hệ thống.
  \item \textbf{Tổ chức học liệu:} nhóm từ vựng thành các bộ thẻ (deck)
        và chủ đề; hỗ trợ chia sẻ, sao chép bộ thẻ trong cộng đồng.
  \item \textbf{Học theo cơ chế ôn tập ngắt quãng:} lập lịch ôn tập tự
        động theo mức độ ghi nhớ từng từ, với nhiều dạng câu hỏi luyện tập
        khác nhau trong một phiên học.
  \item \textbf{Luyện tập chủ động có phản hồi:} luyện đặt câu với từ vựng
        và luyện phát âm, đều có cơ chế chấm điểm để người học nhận phản
        hồi định lượng.
  \item \textbf{Hỗ trợ tạo nội dung bằng AI:} sinh nhanh dữ liệu cho một
        từ chỉ từ dạng nguyên thể, nhập hàng loạt và làm giàu dữ liệu từ
        điển, sinh âm thanh phát âm tự động, nhằm giảm công sức tạo nội
        dung thủ công.
  \item \textbf{Theo dõi tiến độ và tạo động lực:} thống kê số từ theo
        trạng thái học, chuỗi ngày học, biểu đồ hoạt động và bảng xếp hạng
        cộng đồng.
\end{enumerate}

\subsection{Nhu cầu phi chức năng}

\begin{itemize}
  \item \textbf{Bảo mật:} bảo vệ tài khoản và dữ liệu người dùng bằng cơ
        chế xác thực bằng token, phân quyền theo vai trò và kiểm soát truy
        cập theo quyền sở hữu dữ liệu.
  \item \textbf{Hiệu năng và khả năng mở rộng:} đáp ứng tốt khi số lượng
        người dùng và lượng từ vựng tăng; các tác vụ nặng (sinh nội dung,
        sinh âm thanh) được xử lý bất đồng bộ để không chặn luồng chính.
  \item \textbf{Tính nhất quán dữ liệu:} đảm bảo toàn vẹn dữ liệu giữa từ
        vựng, bộ thẻ và tiến độ học của người dùng.
  \item \textbf{Khả năng bảo trì và mở rộng:} kiến trúc rõ ràng, API có
        phiên bản hóa để dễ dàng phát triển thêm tính năng về sau.
\end{itemize}

% ---------------------------------------------------------------------
\section{Mục tiêu của hệ thống}
\label{sec:muc-tieu}

\subsection{Mục tiêu tổng quát}

Mục tiêu tổng quát của đề tài là \textbf{xây dựng hệ thống phía máy chủ
(backend) cho ứng dụng học từ vựng tiếng Anh}, cung cấp đầy đủ các dịch vụ
qua giao diện lập trình ứng dụng (API) để kết hợp giữa kho từ vựng được
biên tập, khả năng cá nhân hóa, cơ chế ôn tập ngắt quãng có cơ sở khoa
học, hỗ trợ sinh nội dung bằng AI và các hình thức luyện tập chủ động có
chấm điểm, qua đó giúp người học ghi nhớ và vận dụng từ vựng hiệu quả hơn.

\subsection{Mục tiêu cụ thể}

\begin{enumerate}
  \item Xây dựng phân hệ \textbf{tài khoản và xác thực} hỗ trợ nhiều
        phương thức đăng nhập, phân quyền người dùng và quản trị viên.
  \item Thiết kế \textbf{mô hình dữ liệu từ vựng} đầy đủ (nhiều nghĩa, bản
        dịch, ví dụ, phiên âm, âm thanh, chủ đề) cho cả kho hệ thống và từ
        vựng cá nhân.
  \item Xây dựng phân hệ \textbf{bộ thẻ và chủ đề}, hỗ trợ chia sẻ và sao
        chép trong cộng đồng.
  \item Cài đặt \textbf{động cơ học theo ôn tập ngắt quãng} dựa trên thuật
        toán SM-2 kết hợp các bước học (learning steps), với nhiều dạng
        câu hỏi luyện tập trong một phiên học.
  \item Xây dựng các phân hệ \textbf{luyện viết câu} và \textbf{luyện phát
        âm} có cơ chế chấm điểm và phản hồi.
  \item Tích hợp các tính năng \textbf{hỗ trợ bằng AI}: tạo nhanh và nhập
        hàng loạt từ vựng, làm giàu dữ liệu từ điển và sinh âm thanh phát
        âm tự động.
  \item Xây dựng phân hệ \textbf{theo dõi tiến độ}: thống kê, chuỗi ngày
        học, biểu đồ hoạt động và bảng xếp hạng.
  \item Đảm bảo các \textbf{yêu cầu phi chức năng} về bảo mật, hiệu năng,
        khả năng mở rộng và tính nhất quán dữ liệu.
\end{enumerate}

\subsection{Phạm vi đề tài}

Đề tài bao gồm cả \textbf{hệ thống phía máy chủ (backend)} cùng các API
phục vụ toàn bộ nghiệp vụ nêu trên và \textbf{ứng dụng giao diện người dùng
phía máy khách (frontend)} tiêu thụ các API đó. Trong đó, trọng tâm trình
bày đặt ở phần backend --- nơi tập trung phần lớn logic nghiệp vụ, mô hình
dữ liệu và các tính năng thông minh; phần frontend được trình bày ở mức kiến
trúc và hiện thực hóa tầng giao diện, đủ để làm rõ cách toàn bộ hệ thống vận
hành từ người dùng cuối đến cơ sở dữ liệu. Các chương tiếp theo sẽ lần lượt
trình bày cơ sở lý thuyết, phân tích và thiết kế hệ thống, cũng như quá
trình hiện thực hóa và kiểm thử.

%  vào main.tex (sau phần mở đầu, trước các chương sau).
%
%  Các nguồn trích dẫn (\cite) trong chương nằm trong references.bib.
%  KHÔNG hardcode số chương/mục/bảng — luôn dùng \label + \ref / \cite.
% =====================================================================

\chapter{TỔNG QUAN ĐỀ TÀI}
\label{chuong:tong-quan}

Chương này trình bày bối cảnh và lý do hình thành đề tài. Trước hết,
mục~\ref{sec:dat-van-de} đặt ra bài toán học từ vựng tiếng Anh và những
khó khăn cốt lõi của người học. Tiếp theo, mục~\ref{sec:khao-sat-hien-trang}
khảo sát hiện trạng các phương pháp học từ vựng đang được sử dụng, còn
mục~\ref{sec:he-thong-lien-quan} phân tích một số hệ thống tiêu biểu trên
thị trường. Trên cơ sở đó, mục~\ref{sec:phan-tich-nhu-cau} phân tích nhu
cầu thực tế và mục~\ref{sec:muc-tieu} xác định mục tiêu mà hệ thống đề xuất
hướng tới.

% ---------------------------------------------------------------------
\section{Đặt vấn đề}
\label{sec:dat-van-de}

Trong quá trình học một ngoại ngữ, từ vựng đóng vai trò nền tảng. Người
học có thể giao tiếp ở mức cơ bản khi nắm vững từ vựng dù ngữ pháp chưa
hoàn chỉnh, nhưng rất khó diễn đạt ý tưởng nếu vốn từ nghèo nàn. Đối với
tiếng Anh — ngôn ngữ phổ biến nhất trong học thuật, công việc và Internet
— việc tích lũy và ghi nhớ một vốn từ vựng đủ rộng là yêu cầu thiết yếu
đối với học sinh, sinh viên và người đi làm tại Việt Nam.

Tuy nhiên, ghi nhớ từ vựng lâu dài là một thách thức. Theo nghiên cứu kinh
điển của Ebbinghaus về \textit{đường cong lãng quên} (forgetting curve),
phần lớn thông tin mới sẽ bị quên đi nhanh chóng nếu không được ôn tập lại
một cách có hệ thống \cite{ebbinghaus1913memory}. Các nghiên cứu về tâm lý
học nhận thức sau này cũng chỉ ra rằng việc \textit{ôn tập ngắt quãng}
(spaced repetition) — phân bố các lần ôn tập theo những khoảng thời gian
tăng dần thay vì học dồn — giúp cải thiện đáng kể khả năng ghi nhớ dài hạn
so với cách học nhồi nhét \cite{cepeda2006distributed}. Đây chính là cơ sở
khoa học cho các thuật toán lập lịch ôn tập được ứng dụng trong nhiều phần
mềm học tập hiện đại.

Mặc dù vậy, cách học từ vựng phổ biến hiện nay của nhiều người Việt vẫn
còn nhiều hạn chế:

\begin{itemize}
  \item \textbf{Học thụ động và rời rạc:} ghi từ ra sổ tay hoặc tra từ
        điển một lần rồi không có cơ chế nhắc ôn tập, dẫn đến nhanh quên.
  \item \textbf{Không cá nhân hóa:} mọi từ được ôn với tần suất như nhau,
        không phân biệt từ người học đã thuộc với từ còn yếu, gây lãng phí
        thời gian.
  \item \textbf{Thiếu luyện tập chủ động:} người học chủ yếu \emph{nhận
        biết} nghĩa của từ (passive) mà ít có cơ hội \emph{sử dụng} từ
        trong câu hoặc luyện phát âm, trong khi đây mới là kỹ năng quyết
        định khả năng vận dụng thực tế.
  \item \textbf{Khó duy trì động lực:} không có công cụ theo dõi tiến độ,
        chuỗi ngày học hay mục tiêu, khiến người học dễ bỏ cuộc giữa chừng.
\end{itemize}

Từ những vấn đề trên, đề tài hướng tới việc \textbf{xây dựng một hệ thống
ứng dụng học từ vựng tiếng Anh} kết hợp giữa cơ chế ôn tập ngắt quãng có
cơ sở khoa học, khả năng cá nhân hóa nội dung học, các hình thức luyện tập
chủ động (viết câu, luyện phát âm có chấm điểm) và công cụ theo dõi tiến độ,
nhằm giúp người học ghi nhớ từ vựng hiệu quả và bền vững hơn.

% ---------------------------------------------------------------------
\section{Khảo sát hiện trạng}
\label{sec:khao-sat-hien-trang}

Để hiểu rõ bài toán, đề tài khảo sát các phương pháp và công cụ học từ
vựng đang được người học sử dụng phổ biến, cùng những ưu, nhược điểm của
chúng.

\subsection{Các phương pháp học từ vựng phổ biến}

\begin{itemize}
  \item \textbf{Sổ tay và giáo trình giấy:} người học chép từ mới kèm
        nghĩa và ví dụ ra giấy. Cách làm này đơn giản, không phụ thuộc
        thiết bị, nhưng không có cơ chế nhắc ôn tập, khó tra cứu và không
        thống kê được tiến độ.
  \item \textbf{Thẻ ghi nhớ giấy (flashcard):} viết từ ở một mặt và nghĩa
        ở mặt còn lại để tự kiểm tra. Hiệu quả hơn sổ tay nhờ tính chủ
        động, nhưng việc sắp xếp thứ tự ôn tập hoàn toàn thủ công và tốn
        công tạo thẻ.
  \item \textbf{Ứng dụng từ điển:} tra nghĩa, phiên âm và phát âm nhanh
        chóng, nhưng bản chất là công cụ tra cứu tức thời, không được
        thiết kế để hỗ trợ ghi nhớ lâu dài.
  \item \textbf{Ứng dụng học ngoại ngữ trên di động:} cung cấp lộ trình
        học, trò chơi hóa và nhắc nhở. Đây là hướng tiếp cận hiện đại
        nhất và là đối tượng được phân tích kỹ ở mục~\ref{sec:he-thong-lien-quan}.
\end{itemize}

\subsection{Những hạn chế chung}

Qua khảo sát, các phương pháp truyền thống và một số công cụ hiện có còn
tồn tại những hạn chế chung sau:

\begin{enumerate}
  \item Phần lớn dựa trên việc \emph{lặp lại máy móc}, thiếu thuật toán
        lập lịch ôn tập tối ưu theo mức độ ghi nhớ của từng từ.
  \item Nội dung học thường \emph{cố định và không thể tùy biến}; người
        học khó tự bổ sung những từ vựng gắn với nhu cầu riêng (chuyên
        ngành, công việc, sở thích).
  \item Việc tạo nội dung học (định nghĩa, ví dụ, phiên âm, âm thanh) tốn
        nhiều \emph{công sức thủ công}, là rào cản lớn khi muốn mở rộng
        kho từ vựng cá nhân.
  \item Các kỹ năng \emph{sản sinh ngôn ngữ} như đặt câu và phát âm hiếm
        khi được luyện tập kèm phản hồi định lượng để người học biết mình
        đang ở đâu.
  \item Thiếu sự gắn kết giữa \emph{theo dõi tiến độ} và \emph{động lực}:
        người học không có chuỗi ngày học, biểu đồ hoạt động hay yếu tố
        cộng đồng để duy trì thói quen.
\end{enumerate}

Những hạn chế này là cơ sở để đề tài xác định các chức năng cần xây dựng,
được trình bày trong mục~\ref{sec:phan-tich-nhu-cau}.

% ---------------------------------------------------------------------
\section{Các hệ thống liên quan}
\label{sec:he-thong-lien-quan}

Trên thị trường hiện có nhiều ứng dụng hỗ trợ học từ vựng và ngoại ngữ.
Phần này phân tích một số hệ thống tiêu biểu để rút ra ưu, nhược điểm và
khoảng trống mà đề tài có thể khai thác.

\subsection{Duolingo}

Duolingo là một trong những nền tảng học ngoại ngữ miễn phí phổ biến nhất,
nổi bật với cách tiếp cận \textit{trò chơi hóa} (gamification): lộ trình
theo cấp độ, điểm kinh nghiệm, chuỗi ngày học và bảng xếp hạng để duy trì
động lực \cite{vonahn2013duolingo}. Tuy nhiên, Duolingo tập trung vào các
bài học có sẵn theo lộ trình cố định; người học khó tự thêm danh sách từ
vựng riêng và việc đào sâu một từ cụ thể (nhiều nghĩa, ví dụ, ngữ cảnh)
còn hạn chế.

\subsection{Anki}

Anki là phần mềm thẻ ghi nhớ mã nguồn mở, áp dụng thuật toán ôn tập ngắt
quãng SM-2 có nguồn gốc từ dự án SuperMemo \cite{wozniak1990optimization,
anki}. Điểm mạnh của Anki là cơ chế lập lịch ôn tập mạnh mẽ và khả năng
tùy biến thẻ rất cao. Nhược điểm là giao diện và quy trình cấu hình tương
đối phức tạp với người mới; toàn bộ nội dung thẻ phải tự tạo thủ công; và
Anki không tích hợp sẵn cơ chế luyện viết câu hay chấm điểm phát âm.

\subsection{Quizlet}

Quizlet cho phép tạo các bộ thẻ từ vựng và học qua nhiều chế độ (thẻ ghi
nhớ, trắc nghiệm, ghép cặp) cùng tính năng chia sẻ bộ thẻ trong cộng đồng
\cite{quizlet}. Quizlet dễ sử dụng và có kho nội dung do người dùng đóng
góp lớn, nhưng cơ chế ôn tập ngắt quãng không sâu bằng Anki và cũng thiếu
các hình thức luyện tập chủ động có chấm điểm.

\subsection{So sánh và khoảng trống}

Bảng~\ref{tab:so-sanh-he-thong} so sánh các hệ thống trên với hệ thống đề
xuất theo những tiêu chí cốt lõi rút ra từ mục~\ref{sec:khao-sat-hien-trang}.

\begin{table}[h]
  \centering
  \small
  \renewcommand{\arraystretch}{1.3}
  \begin{tabular}{|p{4.1cm}|p{2.0cm}|p{1.7cm}|p{1.9cm}|p{2.3cm}|}
    \hline
    \textbf{Tiêu chí} & \textbf{Duolingo} & \textbf{Anki} & \textbf{Quizlet} & \textbf{Hệ thống đề xuất} \\
    \hline
    Ôn tập ngắt quãng (SRS)          & Hạn chế & Mạnh   & Trung bình & Mạnh \\
    \hline
    Tùy biến / tạo từ vựng cá nhân   & Không   & Có     & Có         & Có \\
    \hline
    Hỗ trợ sinh nội dung bằng AI     & Không   & Không  & Hạn chế    & Có \\
    \hline
    Luyện viết câu có phản hồi       & Không   & Không  & Không      & Có \\
    \hline
    Luyện phát âm có chấm điểm       & Hạn chế & Không  & Không      & Có \\
    \hline
    Theo dõi tiến độ \& gamification & Mạnh    & Cơ bản & Cơ bản     & Có \\
    \hline
  \end{tabular}
  \caption{So sánh các hệ thống học từ vựng tiêu biểu với hệ thống đề xuất}
  \label{tab:so-sanh-he-thong}
\end{table}

Như Bảng~\ref{tab:so-sanh-he-thong} cho thấy, mỗi hệ thống có thế mạnh
riêng nhưng chưa hệ thống nào đồng thời đáp ứng tốt cả bốn yếu tố: kho từ
vựng vừa biên tập sẵn vừa cho phép cá nhân hóa, cơ chế ôn tập ngắt quãng
mạnh, hỗ trợ sinh nội dung tự động bằng AI, và tích hợp luyện viết câu cùng
luyện phát âm có chấm điểm. Đây chính là khoảng trống mà đề tài hướng tới.

% ---------------------------------------------------------------------
\section{Phân tích nhu cầu thực tế}
\label{sec:phan-tich-nhu-cau}

Từ hiện trạng và phân tích các hệ thống liên quan, đề tài xác định các nhu
cầu thực tế cần được hệ thống đáp ứng, chia thành nhu cầu chức năng và nhu
cầu phi chức năng.

\subsection{Đối tượng người dùng}

Hệ thống phục vụ hai nhóm người dùng chính:

\begin{itemize}
  \item \textbf{Người học (user):} người dùng cuối, có nhu cầu học và ôn
        tập từ vựng, tạo từ vựng và bộ thẻ cá nhân, luyện tập và theo dõi
        tiến độ của bản thân.
  \item \textbf{Quản trị viên (admin):} người biên tập và quản lý kho từ
        vựng hệ thống (nội dung được kiểm duyệt, dùng chung cho mọi người
        học), quản lý chủ đề và bộ thẻ mẫu.
\end{itemize}

\subsection{Nhu cầu chức năng}

\begin{enumerate}
  \item \textbf{Quản lý tài khoản và xác thực:} đăng ký, đăng nhập bằng
        email/mật khẩu và qua nhà cung cấp bên thứ ba (Google, Apple,
        GitHub); quản lý hồ sơ và thiết lập mục tiêu học tập.
  \item \textbf{Quản lý từ vựng:} mô hình hóa từ vựng phong phú gồm nhiều
        nghĩa, bản dịch, câu ví dụ, phiên âm, âm thanh và gắn chủ đề; cho
        phép người học tự tạo từ vựng cá nhân bên cạnh kho hệ thống.
  \item \textbf{Tổ chức học liệu:} nhóm từ vựng thành các bộ thẻ (deck)
        và chủ đề; hỗ trợ chia sẻ, sao chép bộ thẻ trong cộng đồng.
  \item \textbf{Học theo cơ chế ôn tập ngắt quãng:} lập lịch ôn tập tự
        động theo mức độ ghi nhớ từng từ, với nhiều dạng câu hỏi luyện tập
        khác nhau trong một phiên học.
  \item \textbf{Luyện tập chủ động có phản hồi:} luyện đặt câu với từ vựng
        và luyện phát âm, đều có cơ chế chấm điểm để người học nhận phản
        hồi định lượng.
  \item \textbf{Hỗ trợ tạo nội dung bằng AI:} sinh nhanh dữ liệu cho một
        từ chỉ từ dạng nguyên thể, nhập hàng loạt và làm giàu dữ liệu từ
        điển, sinh âm thanh phát âm tự động, nhằm giảm công sức tạo nội
        dung thủ công.
  \item \textbf{Theo dõi tiến độ và tạo động lực:} thống kê số từ theo
        trạng thái học, chuỗi ngày học, biểu đồ hoạt động và bảng xếp hạng
        cộng đồng.
\end{enumerate}

\subsection{Nhu cầu phi chức năng}

\begin{itemize}
  \item \textbf{Bảo mật:} bảo vệ tài khoản và dữ liệu người dùng bằng cơ
        chế xác thực bằng token, phân quyền theo vai trò và kiểm soát truy
        cập theo quyền sở hữu dữ liệu.
  \item \textbf{Hiệu năng và khả năng mở rộng:} đáp ứng tốt khi số lượng
        người dùng và lượng từ vựng tăng; các tác vụ nặng (sinh nội dung,
        sinh âm thanh) được xử lý bất đồng bộ để không chặn luồng chính.
  \item \textbf{Tính nhất quán dữ liệu:} đảm bảo toàn vẹn dữ liệu giữa từ
        vựng, bộ thẻ và tiến độ học của người dùng.
  \item \textbf{Khả năng bảo trì và mở rộng:} kiến trúc rõ ràng, API có
        phiên bản hóa để dễ dàng phát triển thêm tính năng về sau.
\end{itemize}

% ---------------------------------------------------------------------
\section{Mục tiêu của hệ thống}
\label{sec:muc-tieu}

\subsection{Mục tiêu tổng quát}

Mục tiêu tổng quát của đề tài là \textbf{xây dựng hệ thống phía máy chủ
(backend) cho ứng dụng học từ vựng tiếng Anh}, cung cấp đầy đủ các dịch vụ
qua giao diện lập trình ứng dụng (API) để kết hợp giữa kho từ vựng được
biên tập, khả năng cá nhân hóa, cơ chế ôn tập ngắt quãng có cơ sở khoa
học, hỗ trợ sinh nội dung bằng AI và các hình thức luyện tập chủ động có
chấm điểm, qua đó giúp người học ghi nhớ và vận dụng từ vựng hiệu quả hơn.

\subsection{Mục tiêu cụ thể}

\begin{enumerate}
  \item Xây dựng phân hệ \textbf{tài khoản và xác thực} hỗ trợ nhiều
        phương thức đăng nhập, phân quyền người dùng và quản trị viên.
  \item Thiết kế \textbf{mô hình dữ liệu từ vựng} đầy đủ (nhiều nghĩa, bản
        dịch, ví dụ, phiên âm, âm thanh, chủ đề) cho cả kho hệ thống và từ
        vựng cá nhân.
  \item Xây dựng phân hệ \textbf{bộ thẻ và chủ đề}, hỗ trợ chia sẻ và sao
        chép trong cộng đồng.
  \item Cài đặt \textbf{động cơ học theo ôn tập ngắt quãng} dựa trên thuật
        toán SM-2 kết hợp các bước học (learning steps), với nhiều dạng
        câu hỏi luyện tập trong một phiên học.
  \item Xây dựng các phân hệ \textbf{luyện viết câu} và \textbf{luyện phát
        âm} có cơ chế chấm điểm và phản hồi.
  \item Tích hợp các tính năng \textbf{hỗ trợ bằng AI}: tạo nhanh và nhập
        hàng loạt từ vựng, làm giàu dữ liệu từ điển và sinh âm thanh phát
        âm tự động.
  \item Xây dựng phân hệ \textbf{theo dõi tiến độ}: thống kê, chuỗi ngày
        học, biểu đồ hoạt động và bảng xếp hạng.
  \item Đảm bảo các \textbf{yêu cầu phi chức năng} về bảo mật, hiệu năng,
        khả năng mở rộng và tính nhất quán dữ liệu.
\end{enumerate}

\subsection{Phạm vi đề tài}

Đề tài bao gồm cả \textbf{hệ thống phía máy chủ (backend)} cùng các API
phục vụ toàn bộ nghiệp vụ nêu trên và \textbf{ứng dụng giao diện người dùng
phía máy khách (frontend)} tiêu thụ các API đó. Trong đó, trọng tâm trình
bày đặt ở phần backend --- nơi tập trung phần lớn logic nghiệp vụ, mô hình
dữ liệu và các tính năng thông minh; phần frontend được trình bày ở mức kiến
trúc và hiện thực hóa tầng giao diện, đủ để làm rõ cách toàn bộ hệ thống vận
hành từ người dùng cuối đến cơ sở dữ liệu. Các chương tiếp theo sẽ lần lượt
trình bày cơ sở lý thuyết, phân tích và thiết kế hệ thống, cũng như quá
trình hiện thực hóa và kiểm thử.
